As principais limitações, detalhadas na Seção~\ref{subsec:limitacoes_resultados}, decorrem
do cenário de avaliação. A acurácia foi medida por emulação em computador, com imagens
capturadas de um monitor e apenas 21 fotos de hidrômetros físicos. O firmware medido ainda
executa a localização da janela, e não a região de interesse calibrada adotada como solução.
O enlace foi avaliado em bancada, a 2~m, com um único nó, um receptor de canal único e sem
retransmissão. O consumo foi medido com o nó alimentado pela porta USB e com o módulo de
rádio ligado durante o repouso, de modo que as autonomias são estimativas de um modelo
ideal.

A partir dessas limitações, propõem-se os seguintes trabalhos futuros:

\begin{itemize}
    \item implementar a região de interesse calibrada no firmware do fluxo em borda e medir o
    tempo de leitura, a latência e a carga por leitura, hoje apenas estimados em cerca de
    2~s, 3,3~s e 0,13~mAh no tempo acordado;
    \item suspender o módulo de rádio e os demais componentes da placa durante o repouso e
    medir o consumo com bateria e regulador, para obter a autonomia real dos dois fluxos;
    \item ampliar o conjunto de hidrômetros físicos, com uma instalação fixa, e incluir no
    treino do classificador dígitos claros em fundo escuro e roletes em transição;
    \item avaliar o enlace em campo, a distâncias maiores, com um \textit{gateway}
    multicanal e vários nós, e implementar a retransmissão dos blocos perdidos, medindo o
    seu custo em tempo no ar e em energia;
    \item identificar a causa do tempo de inferência do modelo de localização no
    microcontrolador, de 15,6~s, muito acima do esperado para um modelo desse tamanho;
    \item avaliar o fluxo U-Net + CRNN + CTC, descartado pela memória exigida pela U-Net e
    por um operador ausente no TensorFlow Lite Micro, em um hardware com mais memória ou em
    uma versão mais recente da biblioteca;
    \item implementar a atualização remota dos modelos, que hoje exige regravar o firmware
    de cada nó.
\end{itemize}
